% !PRINTED 293
\weberpage{293}
We now add to the two functions $\varphi, \psi$ a third function $f$, which, set equal to zero, represents a \gesperrt{closed} curve, and we determine the characteristic of the system of functions entirely in the earlier manner, by proceeding along the curve $f$
\weq{1}{k = \tfrac{1}{2}\big[E(f;\varphi,\psi) - A(f;\varphi,\psi)\big] = \tfrac{1}{2}\big[E(f;\psi,\varphi) - A(f;\psi,\varphi)\big].}
Thus, for example, in Fig. 11 the characteristic is $k = 2$.

\weberfigureimage{weber-b1-fig11-12.png}{0.96}

In order to be able to apply our theorems, we must make the curves $\varphi, \psi$ closed in some way by arbitrarily constructed connections, all of which, however, are to run outside $f$. The characteristic of this system is then likewise determined by (1). As an example we complete Fig. 11 by Fig. 12, in which these closing connections are drawn. The points of intersection added in this way may be partly exit points $A(\varphi,\psi)$, partly entry points $E(\varphi,\psi)$. We shall denote their numbers by
\[
A'(\varphi,\psi), \quad E'(\varphi,\psi)
\]
while $A(\varphi,\psi)$ denotes the number of the points of intersection originally present. We must now always have
\weq{2}{A(\varphi,\psi) + A'(\varphi,\psi) = E'(\varphi,\psi)}
We divide the points $A(\varphi,\psi)$ into two groups $A_\alpha(\varphi,\psi)$, $A_\varepsilon(\varphi,\psi)$, of which the first are to lie outside, the others inside $f$. Then we have the following characters of the points of intersection
%CONT>>
